\subsection{1858 Ceza Kanunnamesinin İçeriği}
\indent\indent 1858 Ceza Kanunnamesi, başkanlığını Ahmet Cevdet Paşa'nın yaptığı sekiz üyeli Meclis-i Âli-i Tanzîmat komisyonu tarafından hazırlanmış ve 9 Ağustors 1858 tarihinde yürürlüğe girmiştir. 
Mezkur kanunnamenin kayanğı olarak 1810 Fransız Ceza Kanunu gösterilse de, bu konuya dair sarih bir beyan yoktur. 
Bu çıkarım, iki ceza kanunu arasındaki benzerlikler üzerinden yapılmıştır. 1858 kanunnamesi, suçlar ve cezaların sınıflandırmasınını yaparak önceki kanunnamelerden ayrılmaktadır. Suçlar cinayet,
cünha ve kabahat olmak üzere üç sınıfta toplanmıştır. Cezalar ise idam, kürek, nefy, kalebendlik, hapis, mücâzât-ı nakdiye, tenzil-i rütbe şeklinde sınıflandırılmıştır. Benzer sınıflamayı, 1810 
Fransız Ceza Kanunnamesi'nde de görüyoruz. Bir diğer nokta ise, kanunnamenin belki de en tartışmalı bölümü olan \textit{eşhas hakkında vuku bulan cinayet ve cünhalarla mücâzât-ı müterettibeleri} babının 
ihtiva ettiği maddeler ile Fransız Ceza Kanunu'nun açık benzerliğidir. Ahmet Cevdet Paşa'nın aktardığına göre, kendisi \textit{Avrupa kavânini hükmünce} hükümdara karşı suikast maddesini eklemek istemiş
fakat Şevket Paşa hükümdara suikastın kimsenin aklına gelmemesi gerektiğini söyleyerek buna karşı çıkmıştır. Neticede bu madde kanunnameden çıkarılmıştır. Bir sene sonra vuku bulan Kuleli Vakası'nda bu konu
tekrar gündeme gelmiş, olaya karışanların idam edilebilmesi için kanunnamede madde aranması gerekmiştir. Ahmet Cevdet Paşa, açık bir şekilde Avrupa'dan kanun alınmasına karşı değildir.\footfullcite[220-221]{basak}

Mısır Valisi Mehmet Ali Paşa'nın 1845 yılında çıkardığı ceza kanunnamesinin de bu kanunnameyi etkilemiş olması oldukça olasıdır. \textit{Code de Mehmed Ali} diye anılan bu kanunnamede ırza saldırı, kadı ve şeriattan bahsedilmeden 
kalebendlik ya da kürek ile cezalandırılan bir suça dönüşmüştür. Öte yandan iki kanunname arasındaki benzerlikler, iki kanunnamenin de kaynağının 1810 Fransız Ceza Kanunu olmasından da kaynaklanmış olabilir. İki kanunnamede de açıkça 
görülen husus, şer`i hukuka olabildiğince asgari düzeyde yer verilmesidir.\footcite[223]{basak}

Yukarıda da belirtildiği gibi kanunnamenin en tartışmalı konusu \textit{eşhas hakkında vuku bulan cinayet ve cünhalarla mücâzât-ı müterettibeleri} babında, kişiye karşı işlenen suçlara dair hükümler barındırmasıdır. 
Bu alan şer`i hukukun yetki alanındadır. Bu alanlara dair hükümler eklenerek açık bir şekilde şer`i hukuk alanına müdahele edilmektedir. Bundan dolayı kanunnamenin ilk maddesinde bu konuya açıklık getirilmek istemiştir.
\begin{quote}
Doğrudan doğruya hükümet aleyhine vuku bulan cerâyimin icrây-ı mücâzatı devlete ait olduğu gibi, bir şahıs aleyhinde vuku bulan cerâyimin âsayiş-i umumiyi ihlal eyleresi ciheti dahi kezalik devlete ait olduğundan, tayin ve icrası şer`an emr-i ulül-emre âit olan ta'zirin tayin-i derecâtını dahi işbu Kanunname mütekeffil ve mutazammın olup ancak herhalde şer`an muayyen olan hukuk-u şahsiyeye halel gelmevecektir.\footcite[834]{islam_osmanli_hukuku}
\end{quote}
Burada şahıs aleyhinde işlenen suçların genel asayişi ihlal etmesinden dolayı cezanın tayini ve icrasının devlete ait olduğu belirtilmiştir. Öte yandan tayin ve icranın şer`an belirlenmiş olan şahsi hukuka hiçbir şekilde zeval getirilmeyeceği eklenmiştir. Bu madde farklı şekillerde 
yorumlanabilir. Bir taraftan kanun hazırlayıcılar, şer`i hukuk alanına müdahele etmeye açık bir şekilde cesaret etmektedirler. Öte yandan da şer`i hukukun hiçbir suretle ihlal edilmeyeceğini ekleyerek ilmiye sınıfının tepkisine set çekilmek istenmiştir. 
1858 Ceza Kanunnamesi'nin şeriat hükümlerine atıfları, önceki kanunnamelere göre çok daha azdır. Açıkça görülüyor ki, kalemiye sınıfı ilmiye sınıfının siyaset alanını olabildiğince asgari düzeyde tutmak istemektedir.

Kanunname, mukaddime ve üç bab olmak üzere yirmi sekiz fasıl ve iki yüz atmış dört maddeden oluşmaktadır. Önceki ceza kanunnamelerine göre, daha kapsamlı bir ceza hukuku ihtiva etmektedir.

İkinici fasılda cinayet suçunun cezasının, üçüncü fasılda ise cünha ve kabahat suçlarının cezalarının tafsilatı yapılmıştır. Cezalar idam, kürek, nefy, kalebendlik, hapis, mücâzât-ı nakdiye, ref'i rütbe şeklindedir.
Kürek, kalebendlik ve nefy müebbed ve muvakkat olmak üzere iki çeşittir. Hapis cezası ise asgari yirmi dört saat, azami üç senedir(Madde 34). Bulûğ çağına girmemiş bireyler işledikleri suçun cezasına müstehak değildirler(Madde 40).

Birinci bab, toplum emniyetini ve düzenini bozan suçları ve cezalarını içermekte ve toplam on altı fasıldan oluşmaktadır. Fasıllar aşağıdaki şekildedir:
\begin{enumerate}
    \item Fasıl: Devletin dış güvenliğini ihlal eden cinayet ve cünhalar
    \item Fasıl: Devletin iç güvenliğini ihlal eden cinayet ve cünhalar
    \item Fasıl: Rüşvet
    \item Fasıl: Devlet malının zarar görmesi ve kötüye kullanılması
    \item Fasıl: Memuriyetini ve mekiini suistimal edenler ve memuriyet vazifesini yerine getirmeyenler
    \item Fasıl: Devlet memurlarının ferdlere kötü muamelesi
    \item Fasıl: Devlet memurlarına muhalefet, itaatsizlik ve hakaret
    \item Fasıl: Mahpus kaçırmaya ve cinayet zanlılarını gizlemeye cesaret eden şahıslara dair
    \item Fasıl: Mühür bozma, emanet eşya ve resmi evrağın ahzine cesaret eden şahıslara dair
    \item Fasıl: Yetkisi ve yetkinliği olmadan resmi sıfat kullananlar
    \item Fasıl: Mezhep imtiyazlarına taaruz edenler ile kadim ve muteber eserleri tahrip edenlere dair
    \item Fasıl: Telgraf iletişimini ihlal edenler
    \item Fasıl: Ruhsatsız matbaa açanlar ve matbaalarda zaralı belge basanlar ve mekteplerin talim usullarına dair
    \item Fasıl: Kalpazanlık
    \item Fasıl: Sahtekarlık
    \item Fasıl: Kundakçılık
\end{enumerate}

İkinci bab, şahıslara karşı işlenen cinayet ve cünhalara dair olup, toplam on iki fasıldan oluşmaktadır. Fasıllar aşağıdaki şekildedir:
\begin{enumerate}
    \item Fasıl: Katl, yaralama ve korkutmaya dair 
    \item Fasıl: Çocuk düşürme ile karışık meşrubat ve kefilsiz zehir satanlara dair
    \item Fasıl: Irza geçenlerin cezaları
    \item Fasıl: Hapis ve tutuklama usullarına, çocuk ve kız kaçırma suçlarına dair
    \item Fasıl: Yalan yere şahitlik ve yemin edenlere dair
    \item Fasıl: İftira, küfretme ve sır ifşa edenlere dair
    \item Fasıl: Hırsızlık
    \item Fasıl: İflas ve dolandırıcılıkla suçlananların cezaları
    \item Fasıl: Emniyetin suistimali
    \item Fasıl: Müzayede ve ticarete fesad karıştıranlar
    \item Fasıl: Kumar ve piyango
    \item Fasıl: Mala ve cana zarar verme
\end{enumerate}

Üçüncü bab, toplumun genel huzur ve sağlığına ve zabıtaya muhalefet kabahatlarına dair olup, toplam on bir maddeden oluşmaktadır.

Kanunnamede dikkat çeken unsurların başında cinsel suçların da ırz ve namus emniyetinin kapsamında değerlendirilmesi gelmektedir. Cinsel suçlar, ikinci bab üçüncü fasılda \textit{Hetk-i Irz Edenlerin Mücâzâtı Beyanındadır} başlığı altında tanımlanmıştır.
Önceki kanunnamelerde ırz ve namus kavramı erkekler üzerinden yapılırken, cinsel suçlar ile birlikte kadınlar da gündeme gelmiştir. 1840 Ceza Kanunnamesi'nin üçüncü fasıl birinci maddesinde ırz ve namus ile ilgili madde şu şekildedir:
\begin{quote}
Irz ve namus dahi kişinin canı gibi aziz ve muhterem olarak muhafaza ve vikayesi muktezayı hamiyyet ve mübteğay-ı insaniyetten (vatanperverliğin ve insanlığın gereği) olup, faraza bir adamın itibarına dokunacak söz söylemek ve bir adamı döğmek ve bir şahsa söğmek, onun namusunu hetketmek (namusunu lekelemek) demek olduğundan ve bundan böyle zâbıtan-ı askeriye ve neferât ve kavas ve sair umür-u zaptiye ve raptiye memurları hod be hod (askeri subaylar, erler, kavaslar ve zabıta memurları kendi başına buyruk olarak) kimseyi döğemeyip de kimseye bed lakırdı söyleyemeyip, fakat onların memuriyetleri sokaklarda kavga ve niza vukuunda ve bazı erbab-ı töhmetin zuhurunda (mezunlarının yakalanışından) tutup hiçbir şey yapmaksızın doğruca iktiza eden zâbıt mahalline götürüp teslim etmekten ibaret olacağından fimâ ba'd (bundan sonra) gerek umür-u zaptiye ve raptiyeye memur olan asâkir (askerler) ve kavaslardan ve alel-itlak her hangi rütbe ve sınıfdan olur ise olsun, her kim bir gayrısına söğecek ve itibarına dokunur bir söz söyleyecek olur ise, faraza bu keyfiyet Dersaadet'de vuku bulup da buna mütecâsir olan (cesaret eden) dahi ya Devlet-i Aliyyenin ashab-ı merâtibinden (görevlilerinden) veyahut ulemadan olur ise davası mutlaka Meclis-i Ahkâmr-ı Adliyede görülerek cünhasının cesametine (suçunun büyüklüğüne) ve derecesine göre zâbıtı tarafından beş günden yirmi beş güne kadar mahbus ola. Ve cünhası gayet hafif olduğu tebeyyün ettiği takdirde Meclisce kendisini tekdir ve tevbih ile iktifa oluna.\footcite[812-813]{islam_osmanli_hukuku}
\end{quote}
Bu maddede ırz ve namus, bir adamın itibarı, onuru ve namusu olarak tanımlanmıştır. 1851 Kanun-ı Cedid'in ikinci fasıl birinci maddesinde mezkur maddenin tekrarından öteye gidilmemiş ve  şu şekilde ifade edilmiştir.
\begin{quote}
Kâffe-i tebea-i Devlet-i Aliyye emniyet-i can ve mal ve mahfüziyet-i ırza ve namus hukuk-u şer'iyesine nail olduklarından serbestiyet derecesine varmayarak ber-muktezay-ı bürriyet-i şer'iyye küçük ve büyük herkes kendi hukukunu iddiaya muktedir olup uz ve namus dahi kişinin canı gibi aziz ve muhterem olarak muhafaza ve vikayesi muktezay-ı insaniyetten olmakla, faraza bir adamın itibârına dokunacak söz söylemek ve bila mücip döğmek veya söğmek onun namusunu hetketmek olduğundan o maküle hetk-i namusa mütecasir olan kesânın (kimselerin) cünhası şer'an sabit olup hadd-i şer'i lazım geldiği halde hadd-i şer'isi icra oluna.\footcite[825]{islam_osmanli_hukuku}
\end{quote}
1858 Ceza Kanunnamesi, 1810 Fransız Ceza Kanunu'ndaki \textit{Attentats aux Moeurs} kısmını uyarlanmış hali olan hetk-i ırz faslını içermektedir. Bu oldukça cüretkar bir karardır. Bu kanunnameye kadar zina suçu şer`i hukuka dayanarak mücazat olunurdu. Bu kanunname ile zina suçu mücazatı şer`i mahkemelerden alınmaktır. Burada dikkat çeken unsur ise zina kelimesinin kullanılmamasıdır. 
Zina yerine, hetk-i ırz ve fiil-i şenî kelimeleri kullanılmıştır. 
\begin{quote}
    Madde 198 - Bir adam bir kimseye cebren fiil-i şenî icra eder yani ırzına geçer ise muvakkaten küreğe konulur.\\
    Madde 199 - Böyle cebren fiil-i şenî icrasına tasaddi edip de yed-i ihtiyarında olmayan esbâb-ı mânia haylületiyle fiile çıkmamış olur ise üç aydan ekalli olmamak üzere hapis cezasıyla mücâzât olunur.\\
    Madde 200 - Eğer böyle cebren fiil-i şenî, henüz bir ere tezvic olunmamış kız hakkında vuku bulur ise, buna mütecâsir olan kimse iş bu kürek cezasından başka tazmin vermeye dahi müstehak olur.\\
    Madde 201 — Her kim zükür ve inâsdan genç kimseleri idlal ve iğfal ederek fuhşiyata tahrik ve iğra ve esbab-ı husülünü teshil etmevi itivat ederek âdâb-ı umumiyeye münâfi harekete cesaret eyler ise, bir mahdan bir seneye kadar hapis ile mücâzât olunur. Ve eğer bu suretle idlal ve iğfal kaziyesi baba ya ana veya vasi olan kimselerden zuhur eder ise altı mahdan bir buçuk seneye kadar hapis ile mücâzât kılınır.\footcite[865]{islam_osmanli_hukuku}
\end{quote}
İki yüz birinci maddede açık şekilde zina suçuna dair düzenleme yapılmasına rağmen zina kelimesi kullanılmamıştır. 1860 yılında bu maddeye yapılan ek ile zina kelimesi kullanılmış ve suç artık şer`i husus olmaktan çıkarılmıştır.
\begin{quote}
Zevcesiyle birlikte sâkin olduğu hanede diğer hatun ile zina fiili kabihine melüf olan ve zevcesinin şikâvet-i vakıası üzerine fiil-i mezküru irtikâbı tahakkuk eden zevc, beş mecidiye altınından yüz mecidiye altınına kadar cezay-ı nakdi ahziyle mücâzât olunur.\footcite[866]{islam_osmanli_hukuku}
\end{quote}

Yukarıda da değinildiği üzere, ıskat-ı cenin ve adam öldürme fasılları dışında şer`i hukuka atıf yok denecek kadar azdır. Bu atıflar da ceza hukukun işleyişinden ziyade diyet bedeli, miras paylaşımı gibi hususlardadır.
\begin{quote}
    Madde 171 - Hükm-i kanuni hukuk-u şahsiyeyi iskat edemeyeceğinden maktülün veresesi var ise, onların iddiaları üzerine hukuk-u şahsiye davası muhakeme-i şeriyeye havale olunur.\footcite[861]{islam_osmanli_hukuku}\\
    Madde 181 - Katl ve cerh ve darp fiilleri, bir güne ihtilal ve nehb ve gâret-i emval ile beraber icra olunur ise, bunlara bilhassa cesaret edenlerin görecekleri mücâzâttan başka bu hal-i ihtilalin sebep ve muharrik ve mücibi olanlar dahi ayniyle bu fiilleri işlemiş gibi mücâzât olunurlar. Fakat emr-i kısasta hükmr-i şer'i ne ise icra olunur.\footcite[862]{islam_osmanli_hukuku}\\
    Madde 192 - Bir kimse darp yahut sair güne bir fiil ile bir hâmile hatunun iskat-ı cenin eylemesine sebep olur ise, diyet-i şeriyesi istifa olunduktan sonra, eğer bu teaddisi an kasdin olmuş ise muvakkaten kürreğe konulur.\footcite[864]{islam_osmanli_hukuku}
\end{quote}
